\section{Application to the QED3 curvature}\label{sec:qed}
The lattice input is a probability law of real curvature.
For the all-flux chamber model, the infinite-volume construction
in \cite[Theorem~A(i) and Section~2]{ChamberLattice} gives such a
law at each fixed spacing. Its finite model has antiperiodic
spinors, the chamber fermion operator, and parameters
\[
 q>0,\qquad 0<\rho<1,\qquad w_0,w_1,w_2,w_3>0,\qquad M\in\R.
\]
The photon curvature is $u=dA+(2\pi/q)z$, with all closed
integer fluxes $z$ included. The precise operator and photon
measure are those of the cited section. The parameters may
depend on the spacing but are fixed during each volume limit;
the physical charge scaling is $q_a=e\sqrt a$ for fixed $e>0$.

Every selected real-curvature limit $\nu_a$ from that construction
satisfies
\begin{equation}\label{eq:qed-gaussian-input}
 \E_{\nu_a}\exp\left(\sum_{x,i<j}b_{ij}(x)u_{ij}(x)\right)
 \le\exp\left(\frac12\sum_{x,i<j}|b_{ij}(x)|^2\right)
\end{equation}
for every finitely supported real $b$. The constant is one,
independently of the spacing and the selected volume limit.
For completeness, the passage from finite volume uses only a
fixed finite support: it embeds injectively into sufficiently
large tori, so its counting norm is unchanged. The finite Gaussian
estimate \cite[Proposition~5.1]{ChamberLattice} applies there.
Weak convergence passes the bounded continuous cylinder
function $\min\{R,\exp(\sum b_{ij}u_{ij})\}$ to the limit,
and monotone convergence as $R\to\infty$ gives
\eqref{eq:qed-gaussian-input}.

\subsection{Centered plaquettes and their continuum limits}\label{sec:plaquettes}
An oriented lattice two-cochain has real components
\(u_{\ell m}(x)\), \(x\in\Z^3\), \(\ell<m\), with the usual sign
under reversal of orientation. Its coboundary is
\[
 (du)(x)=u_{23}(x+e_1)-u_{23}(x)
        -u_{13}(x+e_2)+u_{13}(x)
        +u_{12}(x+e_3)-u_{12}(x).
\]
For \(p=(x;\ell m)\) and \(a>0\), set
\begin{equation}\label{eq:plaquette-cells}
 Q_a(p)=a\bigl(x+(e_\ell+e_m)/2+[-1/2,1/2)^3\bigr),\qquad
 f_{[a]}(p)=a^{-3/2}\int_{Q_a(p)}f_{\ell m}(y)\,\dd y.
\end{equation}
The \(\ell m\) component of the step field \(F_a\) equals
\(a^{-3/2}u_{\ell m}(x)\) on \(Q_a(p)\).
The geometric action is the exterior-square action on two-forms.


\begin{corollary}\label{cor:qed-continuum}
Let $a_n>0$ tend to zero. For each $n$, let $\nu_{a_n}$ be
a Borel probability law on real arrays
$(u_{ij}(x))_{x\in\Z^3,\,1\le i<j\le3}$ satisfying
\eqref{eq:qed-gaussian-input}. Let \(F_{a_n}\) be their centered-plaquette step fields as in \eqref{eq:plaquette-cells}.
There are a subsequence and a Borel probability law $\mu$ of
real tempered two-forms on $\R^3$ such that the laws of these
step fields converge weakly to $\mu$ in $H_{-4}$. Here
\[
 \|F\|_{-4}^2=\sum_{i<j}\sum_{k\in\N_0^3}
       (2|k|_1+4)^{-8}|F(h_k\,\dd y_i\wedge\dd y_j)|^2.
\]
Each step-field law and every such limit $\mu$ satisfy
\[
 \E e^{F(f)}\le e^{\|f\|_2^2/2}
 \qquad\bigl(f\in\mathcal S(\R^3;\Lambda^2)\text{ real}\bigr).
\]
Along the selected subsequence, all joint moments and cumulants
of complex Schwartz pairings converge to those of $\mu$.
Their finite-family generating functions and logarithms have
the locally uniform convergence asserted in
\cref{thm:gaussian-convergence}.
In particular, this conclusion applies to any choices of the
infinite-volume real-curvature laws of the chamber model just
described.
\end{corollary}

\begin{proof}
For a finitely supported $b$, let $Y=\sum b_{ij}(x)u_{ij}(x)$
and $v^2=\sum|b_{ij}(x)|^2$. Substituting $tb$ in
\eqref{eq:qed-gaussian-input} gives
$\E e^{tY}\le e^{t^2v^2/2}$ for every real $t$.
As in \cite[Proposition~5.1]{QED3Manuscript}, this implies
\begin{equation}\label{eq:qed-second-input}
 \E|Y|^2\le v^2.
\end{equation}
The one-variable calculation is the first paragraph of the
proof of \cref{prop:subgaussian}.

Apply \cref{prop:step-realization} to each law with $d=r=3$,
$C=1$, the three components $12,13,23$, and their offsets
$(e_i+e_j)/2$. Then apply \cref{prop:compactness} to the sequence
of step fields with $s=3$ and $t=4$, as in the proof of
\cref{thm:lattice-continuum}. This also permits different choices
of law when a spacing occurs more than once.
The resulting weak limit on $H_{-4}$ has exactly the
two-form interpretation and norm stated above.

It remains to retain the exponential estimate. For a real
Schwartz two-form $f$, let $b_{a_n}^{(m)}$ be the finite
cell-coefficient truncations used in
\cref{prop:step-realization}. Their pairings converge almost
surely to $F_{a_n}(f)$. Fatou's lemma and
\eqref{eq:qed-gaussian-input} give
\[
 \E e^{F_{a_n}(f)}
 \le\liminf_{m\to\infty}
        \exp\bigl(\|b_{a_n}^{(m)}\|_{\ell^2}^2/2\bigr)
 =\exp\bigl(\|b_{a_n}\|_{\ell^2}^2/2\bigr)
 \le e^{\|f\|_2^2/2}.
\]
This is the infinite-lattice test extension of
\cite[Section~6.1]{LatticeLimits}. Apply
\cref{thm:gaussian-convergence} to the weakly convergent
subsequence already chosen. It gives the Gaussian bound for
$\mu$ and all the stated correlation and cumulant limits.
\end{proof}


The following consequence adds the lattice identities and reflection
hypotheses. The preceding corollary supplies weakly convergent
subsequences; the assertion below applies to every selected weak limit.

\begin{corollary}\label{cor:plaquette-limit}
Let \(a_j\downarrow0\) and let \(\nu_j\) be laws of real
two-cochains on \(\Z^3\). Suppose that:
\begin{enumerate}[label=(\roman*),leftmargin=*]
\item \(du=0\) almost surely and
\(\E_{\nu_j}e^{\langle h,u\rangle}\le e^{\|h\|_{\ell^2}^2/2}\)
for every finitely supported real plaquette test \(h\).
\item Each \(\nu_j\) is invariant under integer translations,
the oriented cubic action and \(u\mapsto-u\).
\item Each \(\nu_j\) is reflection positive for
\(x_\ell=0\) and \(x_\ell\pm x_m=0\), \(\ell<m\),
on character sums supported on plaquettes all of whose corners
lie strictly on the chosen positive side.
\item The laws of \(F_{a_j}\) are realized on \(H_{-4}\) and
converge weakly there to \(P\).
\end{enumerate}
Then \(P\) is a translation-invariant law of closed real tempered
two-forms, invariant under the cubic group and field-sign change.
It satisfies \(\E_Pe^{F(f)}\le e^{\|f\|_2^2/2}\) and is invariant
and reflection positive at every plane
\(x_\ell=c\), \(x_\ell\pm x_m=c\), \(c\in\R\).
All fixed finite-family moments and cumulants converge along
the given sequence.
\end{corollary}

\subsection{Pairings and equivariance}
\Cref{prop:step-realization} gives the measurable realization and
\(F_a(f)=\langle f_{[a]},u\rangle\). The cell estimate
\eqref{eq:cell-contraction} and the Fatou argument in
\cref{cor:qed-continuum} give
\begin{equation}\label{eq:step-gaussian}
 \E e^{F_a(f)}\le e^{\|f\|_2^2/2}.
\end{equation}
If \(\gamma\) is an affine cubic symmetry of the unit lattice
and \(\gamma_a\) its physical rescaling, define
\((\gamma\cdot h)(p)=h(\gamma^{-1}p)\), where reversing the
orientation of a plaquette reverses the coefficient's sign.
A change of variables gives
\begin{equation}\label{eq:limit-test-equivariance}
 (\gamma_a\cdot f)_{[a]}=\gamma\cdot f_{[a]}.
\end{equation}
Signed coordinate permutations carry cell centers to the corresponding
plaquette centers and carry cubes to cubes, up to null boundaries.
The orientation sign is the exterior-square matrix in both actions.
Integer translations give physical translations in \(a\Z^3\).
These are the exact symmetries used in \cref{sec:stability}.

\subsection{The discrete identity and the reflection support}
Let \(\xi\in C_c^\infty(\R^3;\R)\) denote the three-form
\(\xi\,\dd y_1\wedge\dd y_2\wedge\dd y_3\). Define
\[
 \partial_\ell^a\xi(y)
 =a^{-1}\bigl(\xi(y+ae_\ell/2)-\xi(y-ae_\ell/2)\bigr),\qquad
 (v^a_{23},v^a_{13},v^a_{12})
 =(-\partial_1^a\xi,\partial_2^a\xi,-\partial_3^a\xi).
\]
Then
\begin{equation}\label{eq:limit-Bianchi-error}
 F_a(v^a)=0,\qquad
 \|v^a-d^*\xi\|_2
 \le\frac{a^2}{24}\sum_{\ell=1}^3\|\partial_\ell^3\xi\|_2.
\end{equation}
For the exact identity, assign to each unit cube the number
\[
 \xi_{[a]}(x)=a^{-5/2}\int_{a(x+[0,1)^3)}\xi(y)\,\dd y.
\]
The \(23,13,12\) coefficients of \(v^a_{[a]}\) are
\[
 -(\xi_{[a]}(x)-\xi_{[a]}(x-e_1)),\quad
 \xi_{[a]}(x)-\xi_{[a]}(x-e_2),\quad
 -(\xi_{[a]}(x)-\xi_{[a]}(x-e_3)).
\]
Their finite pairing with \(u\) regroups to
\(\sum_x\xi_{[a]}(x)(du)(x)=0\).
The Fourier multiplier of
\(\partial_\ell^a-\partial_\ell\) is
\(i(2a^{-1}\sin(ap_\ell/2)-p_\ell)\), whose modulus is at most
\(a^2|p_\ell|^3/24\). Plancherel gives the second assertion.

For a mirror with defining function
\(\ell(y)=y_j\) or \(y_j\pm y_k\), take compact tests supported
in \(\{\ell\ge\delta\}\), \(\delta>0\).
A nonzero cell coefficient comes from a cell meeting this set.
For every corner \(q\) of its physical plaquette and every \(y\)
in that cell, \(|\ell(q)-\ell(y)|\le2a\):
each coordinate displacement is at most \(a\), and the normal
has at most two nonzero coordinates of modulus one.
Thus \(2a<\delta\) puts all contributing plaquette corners
strictly on the positive side. There are only finitely many such
plaquettes. The lattice reflection inequality and
\eqref{eq:limit-test-equivariance} give \eqref{eq:rp-approx}.

\begin{proof}[Proof of \cref{cor:plaquette-limit}]
Equation \eqref{eq:step-gaussian} gives the common Gaussian input.
Equation \eqref{eq:limit-Bianchi-error} verifies the identity
hypothesis with \(D^*=d^*\). Equivariance gives cubic symmetry and
mesh translations converging to every continuum translation.
The support check gives the eventual character inequalities at the
nine origin mirrors. \Cref{thm:structural-limit} applies.
Field-sign invariance passes directly through its
characteristic-function identity.
\end{proof}

The real-curvature laws supplied by \cite[Theorem~A]{ChamberLattice}
satisfy the additional hypotheses of \cref{cor:plaquette-limit}.
Together with \cref{cor:qed-continuum}, this gives subsequential
continuum laws with all the stated structural properties. The
reflection passage is the origin-height argument of
\cite[Proposition~6.1]{QED3Manuscript}; translations supply every
height. The conditional corollary imposes no arithmetic restriction
on the spacing sequence.

\subsection{Recovery of a prescribed volume limit}
\label{subsec:fixed-spacing-recovery}\label{sec:recovery}
An infinite-lattice law gives a step-field law by
\cref{prop:step-realization}. To treat a prescribed volume limit
of the finite-torus fields, we also need the converse
identification. We use the bump-function argument of
\cite[Section~6.4]{QED3Manuscript}.

Fix a spacing $a>0$. Let $\mathcal P_2$ be the positively oriented
plaquettes of $\Z^3$, and put
\[
 \Omega=\{u\in\R^{\mathcal P_2}:du=0\},
\]
with its coordinate product topology. For each integer $n\ge4$,
let $\lambda_{a,n}$ be a probability law of real closed
two-cochains on $(\Z/n\Z)^3$ satisfying
\begin{equation}\label{eq:recovery-finite-gaussian}
 \E_{\lambda_{a,n}}e^{\langle\beta,u\rangle}
 \le e^{\|\beta\|_{\ell^2((\Z/n\Z)^3)}^2/2}
 \qquad(\beta\text{ a real torus two-cochain}).
\end{equation}
Both the pairing and the norm sum over positively oriented torus
plaquettes. Denote the law of the periodic extension by
$\nu_{a,n}$. Use the centered interpolation and coefficients
$f_{[a]}$ of \eqref{eq:plaquette-cells}, and write $P_{a,n}$
for its law on $H_{-4}$.
For general Schwartz tests these finite-torus laws obey the
periodized bound \eqref{eq:limit-cutoff-G} below, with physical
period $R=an$. The following identification gives the sharp
$L^2$ bound after the volume limit.

\begin{proposition}\label{prop:fixed-spacing-recovery}
Suppose $n_j\to\infty$ and $P_{a,n_j}\Longrightarrow\mu_a$
weakly on $H_{-4}$. Then there is a probability law $\nu_a$
on $\Omega$ such that
\begin{equation}\label{eq:limit-volume-identification}
 \nu_{a,n_j}\Longrightarrow\nu_a,\qquad
 \mu_a=(\sharp_a)_*\nu_a,
\end{equation}
where $\sharp_a$ is the centered step map, defined as a Borel
map into $H_{-4}$ on a set of full $\nu_a$ measure and set
equal to zero elsewhere. The convergence of cochain laws is
in the product topology and holds along the given sequence.
For every finitely supported real cochain $\beta$,
\begin{equation}\label{eq:limit-lattice-G}
 \E_{\nu_a}e^{\langle\beta,u\rangle}
 \le e^{\|\beta\|_{\ell^2(\mathcal P_2)}^2/2}.
\end{equation}
In particular, $\mu_a$ satisfies the Gaussian bound
$\E_{\mu_a}e^{F(f)}\le e^{\|f\|_2^2/2}$ for real Schwartz
two-forms $f$.
\end{proposition}

\begin{proof}
Discard finitely many terms so that $an_j\ge1$.
The finite-volume step maps take values in $H_{-3}$.
Indeed, the periodization estimate below and the Hermite bound
$p_1(h_ke_\alpha)\le31\lambda_k$ from
\cite[Section~3]{QED3Integration} give
$|F_a(h_ke_\alpha)|\le31\lambda_k\|u\|_{\ell^2}$.
Summing with weight $\lambda_k^{-6}$ proves continuity of
each finite-dimensional step map into $H_{-3}$, since
$\sum_k\lambda_k^{-4}<\infty$.

We use the realization already proved in
\cref{prop:step-realization}. Let $\Omega_3$ be its weighted
square-summability domain with exponent $p=3$, intersected with
$\Omega$. For every $u\in\Omega_3$, the step distribution has the
absolutely convergent pairings \eqref{eq:absolute-pairing}.
Let $\Omega^{(a)}\subset\Omega_3$ be the set on which it has
finite $H_{-3}$ norm, and set $\sharp_a=0$ off this set.
The measurability argument in \cref{prop:compactness} makes
$\Omega^{(a)}$ measurable and $\sharp_a$ a Borel map into $H_{-4}$.
Every periodic cochain lies in $\Omega_3$ and, by the preceding
finite-volume estimate, in $\Omega^{(a)}$.

For each plaquette $\sigma$, choose a real smooth bump $b^\sigma$
with only its corresponding two-form component nonzero, supported
inside $Q_a(\sigma)$ and of integral $a^{3/2}$. Then
$(b^\sigma)_{[a]}=\mathbf1_\sigma$. The map
\begin{equation}\label{eq:limit-step-inverse}
 \mathcal J_a:H_{-4}\longrightarrow\R^{\mathcal P_2},
 \qquad \mathcal J_a(F)(\sigma)=F(b^\sigma),
\end{equation}
is continuous for the product topology, and
$\mathcal J_a\sharp_a u=u$ on $\Omega^{(a)}$.
At finite volume the identities are
\[
 P_{a,n_j}=(\sharp_a)_*\nu_{a,n_j},\qquad
 \nu_{a,n_j}=(\mathcal J_a)_*P_{a,n_j}.
\]
The mapping theorem therefore gives
$\nu_{a,n_j}\Longrightarrow\nu_a:=(\mathcal J_a)_*\mu_a$.
The closed conditions $du=0$ pass to this product limit.

Fix a finitely supported real $\beta$ on the covering lattice.
For all sufficiently large $n_j$, its support embeds injectively
in the torus. Its folded norm then equals its covering
$\ell^2$ norm. Apply \eqref{eq:recovery-finite-gaussian}, pass
the bounded continuous cylinder
$\min\{R,e^{\langle\beta,u\rangle}\}$ to the limit, and let
$R\to\infty$. This proves \eqref{eq:limit-lattice-G}.

The Gaussian bound implies the finite-coefficient second-moment
bound, so \cref{prop:step-realization} gives $\nu_a(\Omega_3)=1$.
The step distribution on this domain satisfies the common
$L^2$ test estimate. Applying \cref{prop:compactness} with $s=3$
therefore gives $\nu_a(\Omega^{(a)})=1$.
The finite-coefficient Fatou argument in
\cref{cor:qed-continuum} now gives
\[
 \E_{\nu_a}e^{(\sharp_a u)(f)}
 \le e^{\|f_{[a]}\|_{\ell^2}^2/2}
 \le e^{\|f\|_2^2/2}.
\]

It remains to prove the reverse equality in
\eqref{eq:limit-volume-identification}. Choose a countable
Schwartz-dense family $\mathcal D_0$ of real compact smooth
two-forms. For $f\in\mathcal D_0$, the test
\[
 f'=f-\sum_\sigma f_{[a]}(\sigma)b^\sigma
\]
is a finite sum and satisfies $f'_{[a]}=0$. Every finite-volume
step field therefore annihilates $f'$. The hyperplane
$\{F:F(f')=0\}$ is closed in $H_{-4}$ and has full $\mu_a$
measure. Intersect these events for $f\in\mathcal D_0$ with
the full-measure event
$\mathcal J_a(F)\in\Omega^{(a)}$. On that single event,
\[
 F(f)=\sum_\sigma f_{[a]}(\sigma)F(b^\sigma)
      =(\sharp_a\mathcal J_aF)(f)
 \qquad(f\in\mathcal D_0).
\]
Both sides are tempered distributions, so density gives
$F=\sharp_a\mathcal J_aF$ there. Taking laws proves the
pushforward identity and the Gaussian assertion for $\mu_a$.
\end{proof}

\subsection{Finite-torus estimates and joint cutoff sequences}
\label{subsec:joint-cutoff-sequences}
The exact $L^2$ Gaussian bound in
\cref{thm:gaussian-convergence} applies to the recovered
infinite-volume laws. Before the volume limit, the Schwartz
test is periodized. We recall the estimate from
\cite[Sections~6.1--6.2]{QED3Manuscript} that controls this step.

Put $R=an$ and $f^R(y)=\sum_{k\in\Z^3}f(y+Rk)$.
The finite-torus coefficient is
\begin{equation}\label{eq:folded-cell-coefficient}
 b_{a,n}([\sigma])
 =\sum_{k\in\Z^3}f_{[a]}(\sigma+nk)
 =a^{-3/2}\int_{Q_a(\sigma)}f^R_{\ell m}(y)\,\dd y
 \quad(\sigma=(x;\ell m)).
\end{equation}
The sum converges absolutely. It gives
$F_a(f)=\langle b_{a,n},u\rangle$, and cell Cauchy--Schwarz
gives $\|b_{a,n}\|_{\ell^2}\le
\|f^R\|_{L^2(\R^3/R\Z^3)}$. With $Q_R=[-R/2,R/2)^3$, define
\begin{equation}\label{eq:limit-periodic-seminorm}
 p_R(f)=\|f\|_2+
        26R^{-2}\||y|^2f\|_{L^2(\R^3\setminus Q_R)}.
\end{equation}
Then
\begin{equation}\label{eq:limit-cutoff-G}
 \E_{P_{a,n}}e^{F(f)}
 \le e^{\|b_{a,n}\|_{\ell^2}^2/2}
 \le e^{p_R(f)^2/2}
 \qquad(f\text{ real}).
\end{equation}
To see the last estimate, put
$c_k=\|f\|_{L^2(Q_R+Rk)}$. Minkowski gives
$\|f^R\|_{L^2(\R^3/R\Z^3)}\le\sum_kc_k$. For $k\ne0$,
$|y|\ge R|k|_\infty/2$ on $Q_R+Rk$, and weighted
Cauchy--Schwarz gives
\[
 \sum_{k\ne0}c_k
 \le4R^{-2}
   \left(\sum_{k\ne0}|k|_\infty^{-4}\right)^{1/2}
   \||y|^2f\|_{L^2(\R^3\setminus Q_R)}.
\]
The shell count $24j^2+2$ gives
$\sum_{k\ne0}|k|_\infty^{-4}=4\pi^2+\pi^4/45$,
whose square root times four is less than $26$.
For $R\ge1$, one has $p_R(f)\le p_1(f)$, and
$p_R(f)\to\|f\|_2$ as $R\to\infty$. The seminorm $p_1$
is continuous on Schwartz space.

\begin{corollary}\label{cor:prescribed-volume-first}
At each spacing $a_i\downarrow0$, take the chamber torus laws
with parameters fixed during that volume limit. Let $\mu_i$
be any prescribed weak $H_{-4}$ volume limit of their
centered step fields. Suppose $\mu_i\Longrightarrow\mu$ on
$H_{-4}$. Then $\mu$ has the conclusions of
\cref{cor:plaquette-limit}, including convergence of all
finite-family moments and cumulants from $\mu_i$ to $\mu$.
Its tensor Schwinger functions satisfy \cref{prop:tensor-growth}.
There are torus sides $n_i\to\infty$ such that
$a_in_i\to\infty$,
\[
 P_{a_i,n_i}\Longrightarrow\mu,
\]
so $\mu$ is also realized by a joint sequence of finite
cutoff laws.
\end{corollary}

\begin{proof}
Apply \cref{prop:fixed-spacing-recovery} at each $a_i$.
It identifies the prescribed law $\mu_i$ with the step image
of a weak cochain limit. The chamber theorem
\cite[Theorem~A]{ChamberLattice} applies to every such limit
with the parameters fixed: it gives closedness, Gaussian
domination, the lattice symmetries, and reflection positivity
at the nine mirror families. Thus these recovered laws meet
all the hypotheses of \cref{cor:plaquette-limit}; that
corollary and \cref{prop:tensor-growth} give the assertions
about $\mu$ and the sequence $\mu_i$.

Fix a metric for weak convergence of probability laws on
$H_{-4}$. Along the defining volume sequence for $\mu_i$,
choose $n_i$ so large that $a_in_i\ge i$ and the distance
between $P_{a_i,n_i}$ and $\mu_i$ is less than $1/i$.
Then $P_{a_i,n_i}\Longrightarrow\mu$ and
$R_i:=a_in_i\to\infty$. This is the diagonal selection in
\cite[proof of Proposition~6.1]{QED3Manuscript}.
\end{proof}

